\section*{Capítulo \ref{ch:b3:electronic-spectroscopy} --- Espectroscopia electrónica y fotofísica}

\begin{solution}{exo:b3:electronic-spectroscopy:1}
(a) Prohibida por la regla de espín ($\Delta S = 1$). (b) Prohibida por la \omterm{thm:b3:electronic-spectroscopy:laporte}{regla de Laporte}
($g \to g$); se observa débilmente gracias al acoplamiento vibrónico. (c) Permitida: $B_{1u}$ es la
representación de $z$ en $D_{2h}$. (d) Prohibida por simetría: $A_2$ no es ninguna de $x$,
$y$, $z$ en $C_{2v}$; la banda es débil.
\end{solution}

\begin{solution}{exo:b3:electronic-spectroscopy:2}
$10^7/349 - 10^7/450 = 28653 - 22222 = \qty{6431}{cm^{-1}}$, \qty{0.80}{eV}.
\end{solution}

\begin{solution}{exo:b3:electronic-spectroscopy:3}
$A = 5700 \times 1 \times 1.0\times10^{-5} = 0.057$; fracción absorbida: $1 - 10^{-0.057} =
0.12$.
\end{solution}

\begin{solution}{exo:b3:electronic-spectroscopy:4}
$k_r = \Phi_F/\tau = \qty{1.5e8}{s^{-1}}$, $\tau_r = \qty{6.7}{ns}$; $k_{nr} = (1 - \Phi_F)/\tau
= \qty{1.0e8}{s^{-1}}$.
\end{solution}

\begin{solution}{exo:b3:electronic-spectroscopy:5}
$\int\varepsilon\,\dd\tilde\nu = 1.0645 \times 5700 \times 5000 = \num{3.03e7}$, luego $f \approx
4.32\times10^{-9} \times 3.03\times10^7 = 0.13$: una transición permitida de intensidad
moderada.
\end{solution}

\begin{solution}{exo:b3:electronic-spectroscopy:6}
El factor de Poisson $\eu^{-S}S^v/v!$ es máximo en $v = \lfloor S\rfloor$. $S = 0.3$: la propia
raya 0--0. $S = 1.5$: $v' = 1$, 1.5 veces la raya 0--0. $S = 5$: $v' = 4$
y 5 por igual, 26 veces la raya 0--0.
\end{solution}

\begin{solution}{exo:b3:electronic-spectroscopy:7}
$\sum k = \qty{4.2e8}{s^{-1}}$: $\tau = \qty{2.4}{ns}$, $\Phi_F = 0.24$, $\Phi_{\mathrm{ISC}} = 0.71$
(y $\Phi_{\mathrm{IC}} = 0.05$).
\end{solution}

\begin{solution}{exo:b3:electronic-spectroscopy:8}
$\Phi = 0.546 \times 0.46 \times (0.050/0.040) \times (1.4266/1.333)^2 = 0.36$.
\end{solution}

\begin{solution}{exo:b3:electronic-spectroscopy:9}
$\tau_0/\tau = 1.429$ y 1.860: lineal en $[Q]$, luego es el propio tiempo de vida el que se
acorta: \omterm{def:b3:electronic-spectroscopy:quencher}{extinción dinámica}. $K_{SV} = 43$ L/mol (con ambos puntos), $k_q = 43/8.0
\times 10^{-9} = \qty{5.4e9}{L.mol^{-1}.s^{-1}}$.
\end{solution}

\begin{solution}{exo:b3:electronic-spectroscopy:10}
\omterm{def:b3:electronic-spectroscopy:quencher}{Extinción estática}: las moléculas que siguen emitiendo viven tanto como antes; la mitad de
ellas están unidas, en el estado fundamental, en un complejo no fluorescente. Con una
constante de asociación $K_a$, la fracción libre es $1/(1 + K_a[Q])$ e $I_0/I = 1 +
K_a[Q]$, la misma forma que la de Stern--Volmer pero con $\tau$ invariable.
\end{solution}

\begin{solution}{exo:b3:electronic-spectroscopy:11}
$\tau_r = \tau/\Phi_F = \qty{14}{ns}$; $\Phi_T = 1 - 0.36 = 0.64$. El tiempo de vida medido está
acortado por el \omterm{def:b3:electronic-spectroscopy:radiationless}{cruce entre sistemas} en competencia; $\tau_r$ es el que daría la emisión
por sí sola.
\end{solution}

\begin{solution}{exo:b3:electronic-spectroscopy:12}
$\frac12\langle\alpha\beta - \beta\alpha|\alpha\beta + \beta\alpha\rangle = \frac12(1 + 0 - 0 - 1) = 0$, usando
$\langle\alpha|\alpha\rangle = \langle\beta|\beta\rangle = 1$, $\langle\alpha|\beta\rangle = 0$ para cada electrón.
El \omterm{def:b3:quantum-model-systems:operator}{operador} dipolar no actúa sobre el espín, de modo que el momento de transición contiene este
factor nulo: $T_1 \leftarrow S_0$ no tiene intensidad salvo si el \omterm{def:b3:many-electron-atoms:spin-orbit}{acoplamiento espín--órbita} mezcla
los estados.
\end{solution}

\begin{solution}{pb:b3:electronic-spectroscopy:1}
\textbf{1.} $A = 5700 \times 1 \times 1.0\times10^{-4} = 0.57$.
\textbf{2.} $1 - 10^{-0.57} = 0.73$.
\textbf{3.} Un $\varepsilon$ de unos pocos miles: una transición $\pi\to\pi^*$ permitida, de intensidad
moderada.
\textbf{4.} $1239.8/349 = \qty{3.55}{eV}$.
\textbf{5.} Solo absorbe en el ultravioleta; la luz visible pasa.
\textbf{6.} Para que la luz absorbida siga siendo proporcional a la concentración y evitar la
reabsorción de la luz emitida (efectos de filtro interno).
\textbf{7.} $S_0 \to S_1$ (absorción, seguida de \omterm{def:b3:electronic-spectroscopy:radiationless}{relajación vibracional}); desde $S_1$:
\omterm{def:b3:electronic-spectroscopy:luminescence}{fluorescencia}, \omterm{def:b3:electronic-spectroscopy:radiationless}{conversión interna}, \omterm{def:b3:electronic-spectroscopy:radiationless}{cruce entre sistemas} hacia $T_1$.
\textbf{8.} $k_r = 0.546/19 \times 10^{-9} = \qty{2.87e7}{s^{-1}}$; $\tau_r = \qty{35}{ns}$.
\textbf{9.} $(1 - 0.546)/\tau_0 = \qty{2.39e7}{s^{-1}}$.
\textbf{10.} $28653 - 22222 = \qty{6431}{cm^{-1}}$.
\textbf{11.} La emisión parte de $S_1$ relajado y termina en niveles excitados vibracionalmente
de $S_0$, después de que la molécula y el disolvente se hayan relajado: el fotón tiene menos
energía, aquí la justa para caer en el visible.
\textbf{12.} $\Phi_F \times (349/450) = 0.42$: un 42\,\% de la energía absorbida.
\textbf{13.} Los puntos suben 0.594 por cada \qty{0.010}{mol/L}: una recta que pasa por 1.
\textbf{14.} Mínimos cuadrados con los cinco puntos: pendiente $K_{SV} = \qty{59.4}{L/mol}$,
ordenada en el origen 1.000.
\textbf{15.} $k_q = K_{SV}/\tau_0 = \qty{3.1e9}{L.mol^{-1}.s^{-1}}$.
\textbf{16.} $\tau = \tau_0/(1 + 59.4 \times 0.040) = \qty{5.6}{ns}$.
\textbf{17.} Midiendo los tiempos de vida: en la \omterm{def:b3:electronic-spectroscopy:quencher}{extinción dinámica}, $\tau_0/\tau$ cae sobre la misma recta
que $I_0/I$.
\textbf{18.} El agua tónica no contiene cloruro que la desactive, y es lo bastante ácida para
mantener protonada la quinina, su forma fluorescente.
\textbf{19.} $k_d = 8 \times 8.314 \times 298/(3 \times 0.890\times10^{-3}) =
\qty{7.4e6}{m^3.mol^{-1}.s^{-1}} = \qty{7.4e9}{L.mol^{-1}.s^{-1}}$.
\textbf{20.} $k_q$ es algo menos de la mitad de $k_d$.
\textbf{21.} Alrededor de un 42\,\%.
\textbf{22.} Más lenta: $k_d \propto 1/\eta$, y $k_q$ solo puede disminuir con ella.
\textbf{23.} \textbf{$k_q/k_d = 0.42$}: el cloruro desactiva la quinina a casi la mitad de la
velocidad de los encuentros controlados por la difusión.
\end{solution}
